\section{Conclusion}
\label{sec:conclusion}

This paper established a prospectively specified, source-locked framework for
pre-deployment design of drive-by railway-bridge scour monitoring. The design
asks which RAW/PAA processing and neural components are useful, which pair of
modeled vehicle-response channels should be retained, how accurately one
passage can estimate several pier-specific support-stiffness-loss targets, and
how performance changes when nominal bearing fixity and a local
flexural-rigidity-loss surrogate are added to the task.

The framework's defining safeguards are four explicit scientific blocks
(F40-S, F40-M, L99-S, and L99-M); semantic state grouping and authenticated
common-random-number pairing only where the catalogues actually overlap; a
sealed outer test; Option-C grouped-development out-of-fold adjudication; a
prospectively fixed channel screen; independent selected-pair hyperparameter
optimization in every block; and a disjoint, report-only post-freeze stability
set. The earlier transport/rescue policy is not part of the experiment.
Track-layer damage, wheel polygonization, suspension damage, and alternative
rail-profile phases are likewise deferred rather than silently treated as
active. One generated FRA-class-4 profile realization is shared by every state
and passage.

The physical claim boundary is equally important. The primary label is modeled
vertical support-stiffness loss, not scour depth; bearing fixity is a nominal
design coordinate; the crack is a local $EI$-loss surrogate; and the two
wheelset rows in \texttt{physical8\_v1} are idealized constrained-motion
acceleration proxies, not validated axle-box sensor observations. They remain
diagnostic payload rows and are excluded from learning and channel selection;
the eligible screen contains six singles and 15 pairs. Optional contemporary
challengers enter only after that screen and use exactly the authenticated
primary F40-S pair.

\pending{Summarize only authenticated four-block outcomes: the RAW/PAA
factorial winners, selected physical channel pair, block-local sealed-test
performance and stability, multi-pier localization, bearing leakage, and the
two registered S--M paired comparisons.}

Within those boundaries, the study is designed to provide auditable answers
conditional on the simulated distribution. Field validation remains the
essential next step: the selected configurations must be confronted with
measured crossings, hardware-specific observation models, varied track
profiles, and the simulation-to-field gap that no simulation study can close
from within.
